\subsection{A certified common kernel of the rigid forgetful tests}

\statement{B.56}{Proved}{The rigid, tripod, resolvent and four-point tests have a common kernel}
Let $H_T\cong F_{13}$ be the pure tetrahedral centralizer,
with the ordered parameter and labels of B.18.
For each subset $S$ consisting of the two poles of an
icosahedral $C_5$ and one of its two five-element orbits,
let
\[
 f_S:H_T\longrightarrow P\Gamma_{0,S}
\]
be the homomorphism obtained by forgetting the other punctures.
There is an explicit $h\ne1$ in $H_T$, with nonzero
integral abelianization, such that $f_S(h)=1$ for all twelve
such $S$.  Consequently $h$ also dies under every rigid
five-point and six-point forgetful map of B.13.
In particular none of these seven-point maps is injective,
and none induces an injective map on the full profinite
completion of $H_T$.

Moreover, the intersection of the kernels of these rigid
forgetful maps and all $495$ four-point homomorphisms of
B.18 is nontrivial, both discretely and in the full profinite
completion.  These assertions concern necessary tests for
$\mathcal K$; the displayed elements do not belong to
$\mathcal K$.

The intrinsic geometric versions of all P3 tripod tests
and the three P4 filling tests for the named tetrahedral
group are identically trivial on $H_T$ and on
$\widehat H_T$.  Thus $h$ already gives a nontrivial
intersection for all test families of B.4(a), and the
element $k$ constructed below remains in that intersection
after adjoining all $495$ four-point tests.
Both of these discrete intersections are free groups of
countably infinite rank.  Their corresponding intersections
of completed test kernels are free profinite groups of
countably infinite rank.
Here P4 on $H_T$ means its tetrahedral quotient action
followed by filling; no extension of a general element
of $H_T$ to an automorphism of $\widehat B$ is asserted.

The rigid intersection is the kernel of a homomorphism to $PB_6$
with thirteen certified braid images; its completed analogue
is the kernel of the completed homomorphism.
This gives a terminating exact membership procedure for every
finite source word.  The rigid kernel and the intersection of
all four-point kernels are incomparable, with explicit witnesses
$h$ and $w$, both discretely and after completion.

\begin{proof}
We give the parameter paths, the exact braid words, and
an independently reproducible certification procedure.
Put
\[
 a=\frac{\sqrt3-1}{2},\qquad b=\frac{\sqrt3+1}{2},\qquad
 E=\{0,\infty,1,-1,i,-i\}
       \cup\{\pm a(1\pm i),\ \pm b(1\pm i)\}.
\]
The tetrahedral parameter curve is $S_T=\mathbb P^1\setminus E$.
We use its base point
$u_*=(371+237i)/100$; a path to the icosahedral base point
identifies its fundamental group with $H_T$.
Changing that path conjugates all the elements below
simultaneously and has no effect on any assertion.

The retained subset is
\[
 S_0=\{1,2,3,4,6,9,12\}.
\]
In the cyclic labels of B.18, labels $1,3$ are the poles,
and $\{2,4,6,9,12\}$ is the first pentagon.
For clarity, the twelve rational point functions are
\[
\begin{split}
(p_1,\ldots,p_{12})={}&\bigg(
 u,-u,u^{-1},\frac{u+i}{u-i},-u^{-1},
 -\frac{u+i}{u-i},\frac{u-i}{u+i},-\frac{u-i}{u+i},\\
 &\hspace{9mm}\frac{u+1}{-iu+i},\frac{u+1}{iu-i},
 \frac{u-1}{iu+i},\frac{u-1}{-iu-i}\bigg).
\end{split}
\]
Normalize labels $2,4,6$ to $\infty,0,1$ by
\[
 M_u(z)=\frac{z-p_4(u)}{z-p_2(u)}
             \frac{p_6(u)-p_2(u)}{p_6(u)-p_4(u)}.
\]
The remaining six finite points, in the label order
$(4,6,1,3,9,12)$, trace a braid in $B_6$.
We multiply all their coordinates by
$1+317i/1000$ solely to choose a generic real projection
for reading crossings.  This fixed complex-linear change
does not change the induced mapping class.

For a rational complex number $c$, define $\eta_c$ to
follow the polygon whose successive vertices are
\[
\begin{split}
 u_*,\quad&c+\epsilon(1+i),\quad c+\epsilon(-1+i),\quad
 c-\epsilon(1+i),\\
 &c+\epsilon(1-i),\quad c+\epsilon(1+i),\quad u_*;
 \qquad \epsilon=\frac{3}{250}.
\end{split}
\]
We use $c=-1$, $c=1$, and
$c=c_\alpha=36603(1+i)/100000$.
The three squares surround respectively $-1$, $1$, and
$\alpha=a(1+i)$, each counterclockwise, and contain no
other point of $E$.
The stems avoid $E$ as well.  This follows directly from
the listed algebraic points; the certificate also checks
that the exceptional polynomial
\[
 u(u^4-1)(u^8+14u^4+1)
\]
has no zero on any of the eighteen rational segments.
Write $e_-,e_+,e_\alpha$ for the three based loops and put
\[
 h=e_\alpha e_+e_\alpha^{-1}e_-^{-1}.
 \tag{B.56.1}
\]
Its integral homology class is
$[e_+]-[e_-]\ne0$ in
$H_1(S_T,\mathbb Z)\cong\mathbb Z^{13}$.
This establishes nontriviality in $H_T$ independently of
any choice of a free generating system of meridians.

Here are the three resulting Artin braid words.
A signed integer $j$ denotes $\sigma_{|j|}^{\operatorname{sgn}j}$;
rows in a displayed array are concatenated in order.
With this notation the certified words are
\[
\begin{split}
 \beta_-={}&\left[\begin{array}{l}
 -2,-4,-1,-2,-3,-2,-1,-2,-3,-2,5,2,4,5,\\
 4,3,2,5,4,3,2,1,-3,-2,1,2,3,4,\\
 3,3,-4,-3,-2,-1,2,3,-1,-2,-3,-4,-5,-2,\\
 -3,-4,-5,-4,-2,-5,2,3,2,1,2,3,2,1,4,2
 \end{array}\right],\\[1mm]
 \beta_\alpha={}&\left[\begin{array}{l}
 -2,-4,-1,-2,-3,-1,-2,-1,-3,5,4,5,4,-1,\\
 -2,-3,-4,-5,4,3,4,-5,-4,-3,-2,-1,-5,3,\\
 1,2,1,3,2,1,4,2
 \end{array}\right],\\[1mm]
 \beta_+={}&[-2,-1,-4,-2,-1,4,3,3,-4,1,2,4,1,2].
\end{split}
\]
Exact free reduction in Artin's action on $F_6$ gives
\[
 \beta_\alpha\beta_+\beta_\alpha^{-1}\beta_-^{-1}=1
                 \quad\hbox{in }B_6.
 \tag{B.56.2}
\]
Here the action is faithful; see
\citet[Sections 1.6 and 1.6.1]{DE_GonzalezMeneses2011}.
The identity is consequently stronger than the required
identity in the seven-puncture sphere mapping class group.

We specify how the path calculation is certified, so that
(B.56.2) does not rely on a numerical braid drawing.
On each segment substitute $u=u_0+(u_1-u_0)t$, $0\le t\le1$,
into the six rational functions $M_u(p_j(u))$.
All coefficients lie in $\mathbb Q(i)$.
For each of the fifteen pairs, the real part of the rotated
coordinate difference is a rational polynomial divided by
a strictly positive squared absolute denominator.
We isolate \emph{all} its roots in $(0,1)$ over $\mathbb Q$.
The certificate verifies that every root is simple, all
crossing intervals are disjoint, no crossing occurs at an
endpoint, and the imaginary difference has constant nonzero
sign on each crossing interval.
Its sign and the derivative sign give the crossing sign.
The moving labels are checked to be adjacent at each crossing,
and their final order is checked against the exact endpoint
coordinates.  Polynomial gcd and real-root counts exclude
collisions, poles and exceptional parameter points on every
segment.
Thus the extracted words are exact.

The program \texttt{certify\_rigid7\_monodromy.py} performs
these rational polynomial checks and the exact Artin action.
The data are recorded in
\texttt{rigid7\_exact\_certificate.json}.
They include the polynomial coefficients, isolating intervals,
crossing signs, paths and braid words.
Its actual output is \texttt{certify\_rigid7\_monodromy.out}.
Every check passes.  The earlier numerical search is retained
as a discovery record and is not used as a proof substitute.
We have therefore proved $f_{S_0}(h)=1$.

We next pass to all the rigid subsets.
The twelve subsets consisting of an antipodal pole pair
and one pentagon form the transitive $A_5$-set $A_5/D$.
Indeed, the stabilizer is $D=C_5$: a rotation preserving
such a subset must preserve its uniquely distinguished
pole pair, and the other coset of $D$ in its dihedral
normalizer exchanges the two pentagons.
The exact enumeration in \texttt{check\_rigid7\_labels.py}
also verifies the orbit size and stabilizer directly from
the matrices and the cyclic labels of B.18; its actual output
and all subsets are attached.
The subgroup $T=A_4$ is transitive on this twelve-element
set, because its intersection with every conjugate of $D$
is trivial.  Hence its orbits have size twelve.
For $t\in T$ and $g\in H_T$, centralization gives a
commutative forgetful diagram, or explicitly
\[
 f_{tS_0}(g)=t_*f_{S_0}(g)t_*^{-1},
\]
where $t_*$ identifies the two punctured spheres.
All twelve kernels in $H_T$ are therefore equal.
This proves $f_S(h)=1$ for every rigid seven-point subset.
Each rigid five-point or six-point subset under consideration
is contained in one of them, so the smaller forgetful map
factors through $f_S$ and also kills $h$.
Residual finiteness of the free group $H_T$ shows that
$h$ remains nonidentity in $\widehat H_T$, whereas every
completed forgetful map kills it.  This proves the claimed
failure of completed injectivity.

Finally let $w\ne1$ be the explicit element of B.19 killed
by all $495$ four-point maps.  It is an iterated commutator,
so $[w]=0$ in $H_1(H_T,\mathbb Z)$.
Transport $h$ to the same base point and put
\[
 k=[h,w]=hwh^{-1}w^{-1}.
 \tag{B.56.3}
\]
The rigid-test kernels and the four-point-test kernels are
normal in $H_T$.  Since $h$ lies in all the former and $w$
lies in all the latter, $k$ lies in every one of these kernels.
It is nontrivial: the centralizer of a nonidentity element
of a discrete free group is cyclic.  To recall the proof,
a centralizer preserves the axis of that element in the
free Cayley tree and acts there by integral translations;
this action is injective because the free group acts freely
on vertices.  If $h=r^m$ generates
such a cyclic centralizer up to a nonzero power, then
$[h]=m[r]\ne0$, so $[r]\ne0$ in the torsion-free abelianization.
A commuting $w=r^n$ with $[w]=0$ would have $n=0$, contrary
to $w\ne1$.
Thus $k\ne1$, and residual finiteness again puts a
nonidentity element in the intersection of all the
completed kernels.

We verify the added tripod and filling assertions with their
domains explicit.  Every element of the pure group $H_T$
acts geometrically on the twelve-puncture sphere group.
Quotienting by the normal closure of the nine discarded
peripherals gives the P3 recipe on any three-subset.
The resulting outer action factors through the pure mapping
class group of the three-puncture sphere, which is trivial.
The continuous completed action is therefore trivial as well.
This argument applies to the whole pure geometric group and
does not require an action on the modular extension
$\widehat B$.

For P4, B.30 identifies $F_T=C_\Gamma(T)$ with the pure
mapping class group of the tetrahedral quotient, which has
one ordinary marked point and three named cone points of
orders $2,3,3$.  In particular $H_T\leq F_T$ has a canonical
geometric outer action on
$J=\widehat{\Gamma_T}=C_2\amalg C_3\amalg C_3$.
Filling either named order-three point, or the order-two
point, leaves a sphere with three special points.
The descended mapping class is pure, hence its outer action
on the filled orbifold group is the identity.  The targets
are $\widehat B,\widehat B,C_3\amalg C_3$, respectively,
as in B.53.  Thus all three intrinsic filled actions are
zero on $H_T$ and, by continuity and density, on its full
completion.  On the common $A_5$-centralizing domain these
are the actions of B.53, by its centerless-extension
comparison.  Outside that domain this construction uses
tetrahedral descent rather than a presumed restriction of
an automorphism of $\widehat B$.
Consequently P3 and these three P4 maps impose no additional
kernel condition within $H_T$ or $\widehat H_T$.
This proves nontriviality of the intersection requested in
B.4(a), and of its intersection with all four-point kernels;
the remaining arguments determine their abstract group types.

We can specify its abstract discrete group type.
Let $N_{\mathrm{rig}}$ denote the intersection requested in
B.4(a), with these intrinsic P3 and P4 maps.
The factorization and conjugation arguments above give
\[
 N_{\mathrm{rig}}=\ker(f_{S_0}),\qquad
 N_{\mathrm{all}}=N_{\mathrm{rig}}\cap
                  \bigcap_{|q|=4}\ker(f_q).
 \tag{B.56.4}
\]
Both are nontrivial normal subgroups of $F_{13}$, since
they contain $h$ and $k$, respectively.
The displayed word for $\beta_+$ is
$b\sigma_3^2b^{-1}$, where
$b=[-2,-1,-4,-2,-1,4]$.
Its seven-puncture mapping class is a Dehn twist about an
essential curve separating two punctures from five,
and therefore has infinite order.
Thus $f_{S_0}(H_T)$ is infinite, so both normal subgroups
in (B.56.4) have infinite index.
A nontrivial normal subgroup of infinite index in a
finitely generated free group has infinite rank.
Indeed, if its covering graph had a finite core, choose
a nonempty cyclically reduced word in the subgroup.
Normality makes that same word a reduced closed loop
at every vertex of the cover.  Every vertex would
therefore lie in the finite core, forcing finite index.
The groups are countable, so their free ranks are
exactly countably infinite.

For the completed tests the same factorization gives
the exact description
$N_{\mathrm{rig}}^{\mathrm{test}}=\ker(\widehat f_{S_0})$,
and adjoining the completed four-point maps gives its
intersection $N_{\mathrm{all}}^{\mathrm{test}}$ with their
kernels.  Both are closed normal subgroups of
$\widehat H_T\simeq\widehat F_{13}$ of infinite index:
the essential twist above still has infinite image in
the residually finite group $P\Gamma_{0,7}$ and hence
in its full completion.

Here a precise profinite freeness theorem does apply.
Shusterman's finite-word theorem states that a closed
infinite-index subgroup of a finitely generated free
profinite group is free profinite of rank $\aleph_0$
if it contains a closed normal subgroup containing
a nontrivial finite word in a profinite basis
\cite[Theorem 1.2; Corollary 3.7 in the primary preprint]{DE_Shusterman2016}.
Choose a free basis of the discrete $H_T$; its images
form a profinite basis of $\widehat H_T$.
The elements $h$ and $k$ are nontrivial finite words
in this basis.  Each completed kernel is normal and
contains the appropriate word, so it contains its
closed normal closure.  Every hypothesis of the theorem
holds, proving that
$N_{\mathrm{rig}}^{\mathrm{test}}$ and
$N_{\mathrm{all}}^{\mathrm{test}}$ are free profinite of
countably infinite rank.  This uses the finite-word
hypothesis; normality and infinite index alone would
not justify the conclusion.

The rank assertion uses a convergent profinite basis,
as in the cited theorem.
No equality of this completed kernel with
$\overline{N_{\mathrm{rig}}}$ is being asserted.
That equality would require a further induced-topology
comparison for the image of $f_{S_0}$.
There is also a strict distinction from the full profinite
completion of either discrete kernel.  A countably generated
abstract free group of infinite rank has uncountably many
homomorphisms to $C_2$, by independent choices on its basis.
Its full completion is therefore not metrizable: a compact
metrizable profinite group has only countably many clopen
sets, hence only countably many open subgroups of index two.
Our completed test kernels are closed in the metrizable
$\widehat F_{13}$.  Thus they are not the full profinite
completions of the discrete kernels, despite the two
countable-rank freeness statements.

Neither $h$ nor $k$ centralizes $A_5$: both are nonidentity
pure mapping classes, whereas B.7 gives
$C_P(A_5)=1$.  Residual finiteness of $\Gamma$ also excludes
an actual noncentralizing mapping class from the completed
centralizer.  These words therefore disprove faithfulness
of the test families; they do not furnish an element of
$\mathcal K$.

\paragraph{Thirteen-generator description.}

The following specifies the homomorphism defining the kernel in B.56
on a geometric free basis, and gives an exact membership procedure
for finite words.  It does not identify the completed kernel with
the closure of its discrete counterpart.

\paragraph{A certified geometric basis.}
Retain the base point $u_*=(371+237i)/100$, radius
$\epsilon=3/250$, and square-meridian paths of B.56.
Put
\[
 A=36603/100000,\qquad B=136603/100000.
\]
The thirteen square centers, in their counterclockwise fan order
as seen from $u_*$, are
\[
 \bigl(-B+Bi,\ i,\ B+Bi,\ -A+Ai,\ -1,\ A+Ai,\ 0,
       -A-Ai,\ -B-Bi,\ A-Ai,\ 1,\ -i,\ B-Bi\bigr).
 \tag{B.56.5}
\]
Let $e_j$ follow the straight stem to the northeast corner of
the $j$th square, traverse its boundary counterclockwise, and
return along the same stem.  These are a free basis of
$\pi_1(S_T,u_*)$.

Here is the exact geometric verification, including the
otherwise essential basis condition.  The thirteen closed
squares are disjoint.  Rational segment--rectangle intersection
tests show that each stem meets its own square only at its
northeast corner and misses every other square.  The thirteen
stems have distinct slopes and a common initial point, so their
interiors are disjoint.  Every square contains precisely its
named exceptional parameter point.  For the irrational points,
this follows from
\[
 \frac{17320508075688772}{10^{16}}<\sqrt3<
 \frac{17320508075688773}{10^{16}},
\]
checked by squaring, and the definitions of $a,b$ in B.56.
All these tests are over $\mathbb Q$.  The northeast ray from
$u_*$ to infinity misses the whole construction.  Cutting along
this ray and the embedded petal tree gives the usual geometric
meridian basis of the thirteen-puncture plane.  Its positively
oriented meridian at infinity is
$e_{14}=(e_1\cdots e_{13})^{-1}$.
In particular the previous certified word is literally
$h=e_6e_{11}e_6^{-1}e_5^{-1}$ in this basis.

\paragraph{All thirteen monodromies.}
Retain the set $S_0$ of B.56 and send labels $2,4,6$ to
$\infty,0,1$.  The six finite labeled points define a map
$S_T\to\operatorname{Conf}_6(\mathbb C)$ and therefore a
specified homomorphism
\[
 \theta:F_{13}\longrightarrow PB_6,\qquad e_j\longmapsto\beta_j.
 \tag{B.56.6}
\]
All thirteen words $\beta_j$ are recorded in
\texttt{all13\_braid\_words.json} and
\texttt{all13\_braid\_words.txt}; the sign convention is that
of B.56.  Their reduced lengths and exponent sums are
\[
\begin{array}{c|rrrrrrrrrrrrr}
j&1&2&3&4&5&6&7&8&9&10&11&12&13\\\hline
|\beta_j|&60&58&26&54&58&36&58&82&88&44&14&72&30\\
\operatorname{exp}(\beta_j)&26&-38&26&-4&2&-4&-8&26&-4&26&2&-38&-4
\end{array}
\]
The program \texttt{certify\_all13\_monodromy.py} checks every
crossing on all $78$ rational path segments by the same exact
root-isolation procedure as B.56.  It excludes collisions,
poles and exceptional parameters on every segment, checks all
endpoint orders, and verifies the identity permutation of
each word.  The three previous words for indices $5,6,11$
agree literally with their original certificate.
The complete polynomial and isolating-interval data are in
\texttt{all13\_exact\_certificate.json}; actual stdout is
\texttt{certify\_all13\_monodromy.out}.

\paragraph{The exact preimage and finite-word membership.}
Let $z=\Delta_6^2=(\sigma_1\cdots\sigma_5)^6$.
Capping the disk at its boundary gives
$PB_6/\langle z\rangle=P\Gamma_{0,7}$, with the cap point
the named puncture $2$; see the capping comparison and center
calculation in \citet[Sections 2.2 and 4.3]{DE_GonzalezMeneses2011}.
There is also a useful splitting check.  Forgetting all but the
fixed strands with labels $4,6$ gives
$\ell_{4,6}:PB_6\to PB_2\cong\mathbb Z$.
Those two normalized points are constantly $0,1$, so
$\ell_{4,6}\theta=0$, whereas $\ell_{4,6}(z)=1$.
Consequently the discrete intersection requested in B.4(a) is
the explicitly specified subgroup
\[
                 N_{\mathrm{rig}}=\theta^{-1}(\langle z\rangle)=\ker\theta.
 \tag{B.56.7}
\]
This gives a terminating exact membership test for every finite
word in the thirteen source generators.  Substitute the thirteen
specified braid words.  In general, centrality is tested by
rejecting exponent sums not divisible by $30$ and otherwise
putting $m$ equal to that sum divided by $30$.
With $v=x_1\cdots x_6$, compare the six freely reduced Artin
images to
\[
                       x_j\longmapsto v^m x_jv^{-m}.
\]
Accept precisely if all six agree.  The Artin action is faithful
and the displayed automorphism is the action of $z^m$, so this
criterion is both necessary and sufficient.  For our normalized
source words, $\ell_{4,6}=0$ forces $m=0$, so the implementation
rejects every nonzero exponent sum and otherwise compares
the Artin images to the six generators themselves.
No efficiency bound
is claimed for this free-word expansion.
\texttt{membership\_rigid7.py} implements it and checks the
full-twist convention, the word $h$, and all fourteen peripheral
meridians.  Its actual output is retained.

Completing (B.56.6), let
$Z_{\mathrm{tw}}=\overline{\langle z\rangle}\leq\widehat{PB_6}$.
This central closure is isomorphic to $\widehat{\mathbb Z}$:
the homomorphism $\ell_{4,6}$ restricts to an isomorphism
$\langle z\rangle\to\mathbb Z$, so it induces the full
profinite topology on this cyclic subgroup.  The universal property of completion
gives
\[
 \widehat{PB_6}/Z_{\mathrm{tw}}\cong\widehat{P\Gamma_{0,7}}.
\]
Indeed finite quotients of the right-hand discrete group are
exactly finite quotients of $PB_6$ killing $z$.
Thus the completed test kernel has the exact embedded description
\[
       N_{\mathrm{rig}}^{\mathrm{test}}
                    =\widehat\theta^{-1}(Z_{\mathrm{tw}})=\ker\widehat\theta.
 \tag{B.56.8}
\]
For the last equality, $\widehat\ell_{4,6}$ is an isomorphism
on $Z_{\mathrm{tw}}$ and vanishes on the image of
$\widehat\theta$.  The independent exact program
\texttt{check\_all13\_pair\_linking.py} checks the initial
label order and all fifteen pair linkings of all thirteen words.
It confirms that the $(4,6)$ linking is zero in every case and
that the full twist has linking one.
The thirteen finite words specify this continuous map completely.
The argument uses the central closure, and does not need a
separate assertion about the full center of $\widehat{PB_6}$.
It supplies no new induced-topology comparison for
$\theta(F_{13})$ or its spherical image.

\paragraph{The rigid and four-point kernels are incomparable.}
There is a short exact test detecting $h$ among the four-point
maps.  For $q=\{1,2,3,5\}$, the cross-ratio of B.18 is
\[
                   \psi_q(u)=\left(\frac{u^2-1}{u^2+1}\right)^2.
 \tag{B.56.9}
\]
It is an unramified degree-four covering
\[
 \mathbb P^1\setminus\{0,\infty,\pm1,\pm i\}
                   \longrightarrow\mathbb P^1\setminus\{0,1,\infty\}.
\]
Its zeros and poles are respectively $\pm1,\pm i$, each of
multiplicity two; the preimages of $1$ are $0,\infty$, also
of multiplicity two.  These account for all ramification.
Filling the other eight source punctures kills $e_6$ and
carries $h$ to the difference of the two distinct meridians
at $1,-1$, nonzero already in homology.  The covering induces
an injection on fundamental groups.  Hence
$\psi_{q*}(h)\ne1$.  Residual finiteness of the target
free group shows that this same finite word survives the
completed four-point test.

In the other direction, form the balanced commutator $w$ of
B.19 using the present fourteen meridians.  Its reduced source
word has length $384$.  It dies in every four-point test.
The auxiliary exact program \texttt{check\_all13\_burau.py}
sends $B_6$ to $\operatorname{GL}_6(\mathbb F_{101})$ by the
unreduced Burau matrices at $t=2$: the block at positions
$j,j+1$ is
$\left(\begin{smallmatrix}-1&2\\1&0\end{smallmatrix}\right)$,
with identity elsewhere.  The program checks all defining
braid relations and all inverse matrices.  It obtains
\[
 \rho(\theta(w))=
 \begin{pmatrix}
 77&19&64&20&74&50\\
 47&60&88&59&42&8\\
 55&6&5&12&79&46\\
 8&98&44&86&18&50\\
 50&0&26&19&92&16\\
 5&62&69&96&93&80
 \end{pmatrix}\ne I.
\]
The exponent sum of $\theta(w)$ is zero, so if it were a
central full-twist power it would be the identity.  This
matrix proves $w\notin N_{\mathrm{rig}}$.  In fact the matrix
fails to commute with every displayed braid generator, so
it also survives the finite braid image modulo the central
image of $z$.  This directly witnesses nontriviality in a
finite quotient of the seven-puncture target.
Thus neither the rigid nor the four-point kernel contains
the other, discretely or after completion.  Their intersection
contains the explicit $776$-letter source word $[h,w]$.
The matrix representation is used only to prove a nonidentity;
its faithfulness is not assumed.
\end{proof}
